\section*{\texorpdfstring{\S\CurrentExerciseGroup}{Section \CurrentExerciseGroup}}
\phantomsection
\addcontentsline{toc}{section}{\texorpdfstring{Exercises for \S\CurrentExerciseGroup}{Exercises for Section \CurrentExerciseGroup}}

\begin{enumerate}
\def\labelenumi{\arabic{enumi})}
\tightlist
\item
  Show that if \(f\) and \(g\) are two numerical functions \(\geqslant 0\) , then
\end{enumerate}

\[
\mu^ {\bullet} (\sup (f, g)) + \mu^ {\bullet} (\inf (f, g)) \leqslant \mu^ {\bullet} (f) + \mu^ {\bullet} (g)
\]

(cf.~Ch. IV, §1, Exer. 1).

\begin{enumerate}
\def\labelenumi{\arabic{enumi})}
\setcounter{enumi}{1}
\tightlist
\item
  Show that, for every function \(f \geqslant 0\) defined on \(T\) , the mapping \(\mu \mapsto \mu^{\bullet}(f)\) of \(\mathcal{M}_{+}(T)\) into \(\overline{\mathbf{R}}\) is lower semi-continuous for the quasi-strong topology (Ch. III, §1, Exer. 8; cf.~Ch. IV, §1, Exer. 6 b)). Under what condition is this mapping continuous?
\end{enumerate}

¶ 3) a) Let \((f_{\alpha})_{\alpha \in \mathbb{A}}\) be a family of numerical functions \(\geqslant 0\) , directed for the relation \(\leqslant\) , such that the mapping \(t \mapsto (f_{\alpha}(t))\) of \(\mathbf{T}\) into \(\mathbf{R}^{\mathbf{A}}\) is \(\mu\) -measurable. Let \(f\) be the upper envelope of the family \((f_{\alpha})\) . Then \(f\) is \(\mu\) -measurable and

\[
\int^ {\bullet} f d \mu = \sup _ {\alpha \in A} \int^ {\bullet} f _ {\alpha} d \mu .\tag{1}
\]

\begin{enumerate}
\def\labelenumi{\alph{enumi})}
\setcounter{enumi}{1}
\tightlist
\item
  Let \((f_{\alpha})_{\alpha\in\mathrm{A}}\) be a family of numerical functions \(\geqslant0\) defined on T and having the following property: for every compact subset K of T and every \(\varepsilon>0\) , there exists a compact set \(K'\subset K\) such that \(\mu(K-K')\leqslant\varepsilon\) and such that the restrictions to \(K'\) of all the functions \(f_{\alpha}\) are lower semi-continuous. Let f be the upper envelope of the family \((f_{\alpha})\) . Show that f is measurable and satisfies the relation (1).
\end{enumerate}

Give an example of a mapping \(t \mapsto (f_{\alpha}(t))\) , satisfying the above condition, that is not \(\mu\) -measurable.

\begin{enumerate}
\def\labelenumi{\arabic{enumi})}
\setcounter{enumi}{3}
\tightlist
\item
  Let \(\mathbf{T}\) and \(\mu\) be the locally compact space and measure defined in Exer. 5 of Ch. IV, §1.
\end{enumerate}

\begin{enumerate}
\def\labelenumi{\alph{enumi})}
\item
  Show that \(\mu\) is the supremum of a sequence of measures with finite support, but that \(\mu^{*} \neq \mu^{\bullet}\) , and that \(\mu\) is not moderated. From this, deduce that Prop. 11 becomes false when the symbol \(\int^{\bullet}\) therein is replaced by \(\int^{*}\) , even when the set A is countable.
\item
  For every real number \(y\) , let \(f_y\) be the characteristic function of the set reduced to the point \((0, y)\) of \(\mathbf{T}\) . Show that the mapping \(t \mapsto (f_y(t))_{y \in \mathbf{R}}\) is \(\mu\) -measurable, but that if \(f = \sup_{y \in \mathbf{R}} f_y\) then \(\mu^*(f) = +\infty\) and \(\sup_{y \in \mathbf{R}} \mu^*(f_y) = 0\) .
\end{enumerate}

\begin{enumerate}
\def\labelenumi{\arabic{enumi})}
\setcounter{enumi}{4}
\tightlist
\item
  Let \((\mu_n)\) be a sequence of positive measures on \(T\) that converges quasi-strongly to \(\mu\) in \(\mathcal{M}(T)\) (Ch. III, §1, Exer. 8).
\end{enumerate}

\begin{enumerate}
\def\labelenumi{\alph{enumi})}
\item
  Show that if a mapping \(f\) of \(\mathbf{T}\) into a topological space \(\mathbf{G}\) is \(\mu_n\) -measurable for every \(n\) , then it is \(\mu\) -measurable (make use of Exercise 6 b) of Ch. IV, §1).
\item
  Let \(\mathbf{f}\) be a mapping of T into a Banach space \(F\) , that is essentially \(\mu_n\) -integrable for every \(n\) . Show that if the sequence \((\mu_n^\bullet (|\mathbf{f}|))\) is bounded, then \(\mathbf{f}\) is essentially \(\mu\) -integrable (Exer. 2). Show by means of an example that one does not necessarily have \(\int \mathbf{f} d\mu = \lim_{n\to \infty}\int \mathbf{f} d\mu_n\) (take for \(\mu_n\) a point measure on a compact space, such that \(\lim \| \mu_n\| = 0\) ); however, this relation holds when \(\mathbf{f}\) is bounded and has compact support.
\end{enumerate}

\begin{enumerate}
\def\labelenumi{\arabic{enumi})}
\setcounter{enumi}{5}
\tightlist
\item
  Let \(\mathfrak{K}\) be a \(\mu\) -dense set of compact subsets of T. Let \(\mathbf{f}\) be a function defined on T, with values in a Banach space F, such that \(\mathbf{f}\varphi_{\mathrm{K}}\) is \(\mu\) -integrable for every \(\mathrm{K} \in \mathfrak{K}\) . If the integrals \(\int \mathbf{f}\varphi_{\mathrm{K}} d\mu\) have a limit in F with respect to the directed set \(\mathfrak{K}\) (for \(\subset\) ), then \(\mathbf{f}\) is essentially \(\mu\) -integrable.
\end{enumerate}

\setcounter{exercisegroup}{2}
\edef\CurrentExerciseGroup{\arabic{exercisegroup}}
